\documentclass{article}

%===========================================================
% 導言區 (Preamble)
%===========================================================

\usepackage{xcolor}
\usepackage{tikz}
\usepackage{ctex}

\usetikzlibrary{
    tikzmark,
    positioning,
    arrows.meta
}

% --- 中文字體設定 ---
% Windows: "Microsoft JhengHei UI" / macOS: "PingFang TC"
\setCJKsansfont{PingFang SC}
\setCJKmainfont{PingFang SC}

\pagestyle{empty}

% --- 顏色定義 ---
\definecolor{bgPercent}{rgb}{1.0, 0.8, 0.6}
\definecolor{bgS}{rgb}{0.7, 1.0, 0.7}
\definecolor{bgPattern}{rgb}{0.85, 0.8, 1.0}
\definecolor{bgText}{rgb}{1.0, 1.0, 0.7}
\definecolor{bgG}{rgb}{1.0, 0.7, 0.85}

%===========================================================
% 正文 (Body)
%===========================================================

\begin{document}

\begin{tikzpicture}[remember picture, overlay]
\node at (current page.center) {

% --- Vim 命令 ---
\begin{Huge}
    \setlength{\fboxsep}{2pt}
    \texttt{\tikzmarknode{colon}{:\strut}}%
    \tikzmarknode{percent}{\colorbox{bgPercent}{\texttt{\%\strut}}}%
    \tikzmarknode{s}{\colorbox{bgS}{\texttt{s\strut}}}%
    \tikzmarknode{pattern}{\colorbox{bgPattern}{\texttt{/<pattern>\strut}}}%
    \tikzmarknode{text}{\colorbox{bgText}{\texttt{/<text>\strut}}}%
    \tikzmarknode{g}{\colorbox{bgG}{\texttt{/g\strut}}}%
\end{Huge}

% --- 批注層 ---
\begin{tikzpicture}[remember picture, overlay,
    annot/.style={font=\small\sffamily, text=black!70},
    arr/.style={-Stealth, thick, draw=gray!60}
  ]
    % --- 修改的部分在這裡 ---
    % 透過設定不同的 `below` 數值 (例如 1em 和 2.5em)，
    % 將批注交錯排列在不同的高度，避免重疊。
    \node[annot, below=2.5em of colon]   (annot_colon)   {命令列模式};
    \node[annot, below=1em of percent] (annot_percent) {范围};
    \node[annot, below=2.5em of s]       (annot_s)       {命令};
    \node[annot, below=1em of pattern] (annot_pattern) {目标模式};
    \node[annot, below=2.5em of text]    (annot_text)    {替换内容};
    \node[annot, below=1em of g]       (annot_g)       {global全局};

    % 箭頭的程式碼無需改動，它們會自動延伸到新的位置
    \draw[arr] (annot_percent.north) -- (percent.south);
    \draw[arr] (annot_s.north)       -- (s.south);
    \draw[arr] (annot_pattern.north) -- (pattern.south);
    \draw[arr] (annot_text.north)    -- (text.south);
    \draw[arr] (annot_g.north)       -- (g.south);
    \draw[arr] (annot_colon.north)   -- (colon.south);
\end{tikzpicture}

}; 
\end{tikzpicture} 

\end{document}